
            \documentclass{article}
            \usepackage[utf8]{inputenc}
            \usepackage{hyperref}
            \usepackage{verbatim}

            \begin{document}



d. Press \texttt{Create Datasets}. A list of datasets appears, one line for each set of science data.

\begin{verbatim}
   :alt: create_datasets
   :name: fig_create_datasets
\end{verbatim}

e. Choose the datasets that should be processed 

\begin{verbatim}
   :alt: send
   :name: fig_send
\end{verbatim}
and send them to the data reduction queue by pressing \texttt{Submit to Reduction Queue}. Note that this action does not start the 
reduction automatically.

\subsection{3. Start the reduction}

a. Press \texttt{Reduction Queue}.

\begin{verbatim}
   :alt: processing_queue
   :name: fig_processing_queue
\end{verbatim}


b. (Optional). Configure the workflow and recipe parameters by pressing the wheel 
button ![](figures/configure_dataset.jpg) to open the configuration editor.

c.  Press the \texttt{Reduce} button. The selected data will now be processed with the default parameters.

\begin{verbatim}
   :alt: reduce
   :name: fig_reduce
\end{verbatim}


Congratulations! You reduced your first data with the \texttt{EDPS} dashboard! All the reduced data are saved in the EDPS_data directory specified when executing \texttt{edps-gui} for the first time.


            \end{document}
